\section{True divisors and rational subfields}
\label{sec:good-functions}

The curve-subfield correspondence lets us distinguish valuations trivial
on the constants from the quasi-divisorial valuations recovered in
Section~\ref{sec:local}.  Their restrictions to curve subfields will then
identify a family of rational functions with enough point information to
recover field operations.

A \emph{true prime divisor}, or \emph{divisorial valuation}, of $F/\kappa$
is a quasi-prime divisor trivial on $\kappa$.  Its value group is $\Z$;
we use the discrete normalization when writing its order character.

The true-divisor criterion below has as a methodological antecedent
Topaz's rational-character criterion \cite[Lemma~A.7]{Topaz2021Torelli},
via the companion reconstruction manuscript \cite[Sections~4--6]{InternalBP}.
Here finite coefficients require the finite-dimensional-kernel argument
given below; the rational-coefficient statement is not being invoked as
a proof of this proposition.

\begin{proposition}[Recognition of true divisors]
\label{prop:true-divisors}
Let $F/\kappa$ be one of $K/k$ or $L/l$.  A quasi-divisorial valuation
$v$ of $F/\kappa$ is a true prime divisor if and only if there is a curve
subfield $E\subset F$ for which
\begin{equation}
\label{eq:good-true-test}
                         \dim_\Lambda(V_E\cap D_v^\perp)<\infty.
\end{equation}
Consequently $\Phi$ matches the inertia and decomposition groups of true
prime divisors of $L/l$ with those of $K/k$.
\end{proposition}

\begin{proof}
Suppose first that $v$ is trivial on $\kappa$.  Since the residue field
has transcendence degree $d-1\ge1$, choose a unit $t$ with transcendental
residue, and let $E$ be the relative algebraic closure of $\kappa(t)$
in $F$.  The valuation is trivial on $\kappa(t)$ and therefore on its
algebraic extension $E$.  Reduction embeds $E$ into $Fv$.  Its induced
map $V_E\to V_{Fv}$ has finite-dimensional kernel by
Lemma~\ref{lem:kummer-kernel}.  Lemma~\ref{lem:residue-duality} identifies
that kernel with $V_E\cap D_v^\perp$, proving
\eqref{eq:good-true-test}.

Conversely, suppose that the restriction of $v$ to $\kappa$ is nontrivial,
and let $E$ be any curve subfield.  Choose $x\in E\setminus\kappa$ with
$v(x)\ge0$, replacing a nonconstant element by its inverse if necessary.
There are infinitely many $c\in\kappa^\times$ with $v(c)>0$, for
example the positive powers of one such element.  Each $1+cx$ is a
principal unit for $v$.  Their classes are independent in
$V_{\kappa(x)}$: up to constant factors they are $[x+c^{-1}]$, with
distinct zeros on the rational curve.  The map to $V_E$ has only a
finite-dimensional kernel by Lemma~\ref{lem:kummer-kernel}.  These
principal units therefore span an infinite-dimensional subspace of
$V_E\cap D_v^\perp$.  This disproves \eqref{eq:good-true-test} for every
$E$.  The final assertion now follows from
Theorem~\ref{thm:curve-matching} and the already established matching of
quasi-divisorial pairs.
\end{proof}

For a curve subfield $E\subset F$, restriction of characters gives a
continuous surjection $A_F\to A_E$, dual to the inclusion $V_E\subset V_F$.
If a true divisor restricts nontrivially as a valuation on $E$, that
restriction is a positive integer multiple of a point valuation of the
smooth projective curve of $E$.  Its inertia image in $A_E$ is either
the corresponding point line or zero; the latter occurs when that
integer is divisible by $\ell$.  Thus we retain the nonzero inertia
images when discussing which point lines are detected.

\begin{definition}
\label{def:good-subfield}
A curve subfield $E\subset F$ is \emph{good} if it is rational over
the constants and every point line in $A_E$ is the nonzero image of the
inertia of some true divisor of $F$.  A \emph{good function} is an element
$t\in F$ generating a good subfield $\kappa(t)$.
\end{definition}

\begin{proposition}[Recognition of good subfields]
\label{prop:good-matching}
The curve-subfield correspondence of Theorem~\ref{thm:curve-matching}
restricts to a bijection on good subfields.  For matched good subfields
$E\subset K$ and $P\subset L$, the induced character isomorphism gives
a bijection between the points of their smooth projective rational curves.
\end{proposition}

\begin{proof}
For each curve subfield $E$, let $\mathcal Q_E$ be the collection of
nonzero images of true divisorial inertia in $A_E$.  By
Proposition~\ref{prop:true-divisors}, these collections correspond under
the induced character isomorphisms.  They are subfamilies of the point
lines of $E$.  Choose a nonzero generator of each line.  The generators
tend to zero outside finite sets, so summation defines a continuous map
from a product of copies of $\Lambda$ into $A_E$.

By Lemma~\ref{lem:curve-orders}, when every point occurs, this summation
map has one-dimensional kernel, generated by a coefficient family with
every coordinate nonzero.  A proper subfamily of the point lines has no
relation: extending any proposed relation by zero to the missing points
would contradict the description of the full relation kernel.  Hence
the stated kernel property recognizes exactly when $\mathcal Q_E$
contains every point line.  Once this holds, the common kernel of their
orders on $V_E$ has dimension $2g(E)$.  Its vanishing recognizes genus
zero, equivalently rationality over the algebraically closed constants.
Both conditions are preserved by the character isomorphisms.  Finally,
distinct points have distinct inertia lines, so their matching gives the
asserted point bijection.
\end{proof}

We now construct enough good functions for the later reconstruction.
The construction is local on a smooth model: a ratio of two transverse
parameters has a multiplicity-one divisor over every value, and the
exceptional divisor of the blowup ensures that the resulting rational
subfield is relatively algebraically closed.

\begin{lemma}[A supply of good functions]
\label{lem:good-functions}
Let $F/\kappa$ be a finitely generated extension of an algebraically
closed field with $\trdeg(F/\kappa)\ge2$ and
$\charac(\kappa)\ne\ell$.  On an integral model of $F/\kappa$, let
$q$ be a smooth closed point and let $a,b$ be rational functions regular
at $q$, vanishing there, with independent differentials in the cotangent
space at $q$.  Then $t=a/b$ is a good function.  Moreover, for every
$f\in F^\times$ there is a good function $t$ such that $ft$ is good.
Consequently good functions generate $F$ over $\kappa$, and their classes
span $V_F$.
\end{lemma}

\begin{proof}
Blow up the smooth point $q$ and let $v$ be the valuation of the
exceptional divisor.  Write $e=\trdeg(F/\kappa)$.  Its residue field is
the rational field of $\Pj_\kappa^{e-1}$.  The leading linear forms of
$a$ and $b$ give independent homogeneous coordinates on the exceptional
divisor.  Their ratio is the residue of $t$ and is a rational coordinate
in $Fv$.  In particular $v$ is trivial on $\kappa(t)$, and
$\kappa(\overline t)$ is relatively algebraically closed in $Fv$.
The relative algebraic closure of $\kappa(t)$ in $F$ is a finite
algebraic extension on which $v$ is still trivial.  It therefore embeds
by reduction into $Fv$, over $\kappa(t)\cong\kappa(\overline t)$.
Relative algebraic closedness in the residue field forces this extension
to be $\kappa(t)$ itself.

For each $c\in\kappa$, the function $a-cb$ has nonzero linear term at
$q$ and cuts out a prime divisor locally there, with order one.
The function $b$ is a unit at its generic point because the linear terms
of $a$ and $b$ are independent.  The resulting true divisor has
$\ord(t-c)=1$ and detects the point $c$ of $\Pj^1_\kappa$.
The divisor locally cut out by $b$ similarly detects infinity: $a$
is a unit at its generic point and $t$ has a simple pole.  Every point
line is thus detected, proving that $t$ is good.

Given $f\in F^\times$, choose the smooth point $q$ in an open set where
$f$ is a unit, and choose $a,b$ as above.  The functions $fa,b$ also
vanish at $q$ with independent differentials, since
$d(fa)_q=f(q)\,da_q$.  Hence both $ft=(fa)/b$ and $t=a/b$ are good.
The equality $f=(ft)/t$ proves field generation, and
$[f]=[ft]-[t]$ proves the spanning assertion for $V_F$.
\end{proof}

\begin{lemma}[A common partner]
\label{lem:common-partner}
For any finite list of good functions $t_1,\ldots,t_m$ of $F/\kappa$
there is a good function $s$ such that $s$ is algebraically independent
of each $t_i$ over $\kappa$ and
\[
               t_i+cs\ \hbox{is good for every }i
               \hbox{ and every }c\in\kappa^\times.
\]
\end{lemma}

\begin{proof}
Choose a smooth closed point $q$ on a model where all the $t_i$ are
regular, and choose $a,b$ vanishing at $q$ with independent differentials.
Put $s=a/b$.  It is good by Lemma~\ref{lem:good-functions}.  For
$c\ne0$, write
\[
                             t_i+cs=\frac{t_i b+ca}{b}.
\]
At $q$ the numerator has differential $t_i(q)\,db_q+c\,da_q$, which
is independent of $db_q$.  The same lemma therefore makes every indicated
sum good.

Let $v$ be the exceptional-divisor valuation at $q$ used in the preceding
lemma.  The residue of $s$ is transcendental over $\kappa$.  By contrast,
the restriction of $v$ to $\kappa(t_i)$ is centered at the finite point
$t_i(q)$: the nonzero function $t_i-t_i(q)$ has positive order at $q$.
Its residue field is consequently $\kappa$.  If $s$ were algebraic over
$\kappa(t_i)$, relative algebraic closedness of that good subfield in
$F$ would imply $s\in\kappa(t_i)$, contradicting its transcendental
residue.  This proves the required independence.
\end{proof}
